% Appendix 1: the derivation of the interval likelihood.
%
% Created 2026-08-09 by Phase 1 of papers/1_method/08_length_plan.md. This file holds the algebra
% that used to sit in the body of the Theory section. Nothing here is new text: every block was
% moved verbatim, one commit per move, so that `git log --follow` shows the relocation rather than a
% deletion and an addition. What Theory keeps is the vocabulary the Results are written in (the two
% axes, the eight members, the two closures, the observable) plus a summary of how the algorithms
% were derived; what a reader has to CHECK rather than FOLLOW is here.
%
% ORDER. This appendix is input BEFORE 09_appendix_members so that it is Appendix 1: the
% member-by-member spec cites twelve equations of this derivation and eight of them are cited from
% nowhere else, so with the spec first every one would be a forward reference.
%
% THE \appendix DECLARATION AND THE A-NUMBERING LIVE HERE AND NOWHERE ELSE. They were in
% 08_appendix_members.tex when that was the only appendix; duplicating them in two files restarts
% the counter and produces two equations called A1.
%
% FLOATS. A `figure` environment inside appendixbox is a hard LaTeX error ("Not in outer par mode"),
% because elife.cls wraps the box in mdframed: the caption is dropped and the label goes undefined.
% Anything figure-shaped that lands here must be `center` + \captionof, the pattern
% 09_appendix_members.tex already uses for its table, and must NOT carry `fullwidth`, which escapes
% the panel and overprints the left rule with no overfull warning.

\appendix
\setcounter{equation}{0}
\renewcommand{\theequation}{A\arabic{equation}}

\begin{appendixbox}
\label{app:derivation}
\section*{The interval likelihood, derived}

The Theory section states what each member of the family conditions on and what that buys. This
appendix carries the derivation behind those statements: the boundary-state construction and the
moments it needs, what separates the start-conditioned members from the boundary-conditioned one,
and the correspondence with the linear-filtering literature. It is written to be checked rather
than read in sequence, and every equation the body or Methods refers to is numbered here.

\subsection*{The boundary-conditioned moments, and how they are evaluated}

The ladder needs two conductance objects, both built from the boundary-conditioned
single-channel moments $\bar{\gamma}_{i\to j}$ (the mean interval-averaged
conductance of a channel that starts in $i$ and ends in $j$) and
$\bar{v}_{i\to j}$ (its variance). Both are moments of the integrated conductance
$A_{0\to\Delta}$ of Eq.~\ref{eq:integrated-conductance}, taken conditional on the chain starting
in $i$ and ending in $j$, and both are available in closed form, with no object larger
than $K\times K$ formed anywhere. They are derivatives at $\alpha=0$ of the tilted semigroup
met above. Writing
\begin{equation}
  \mathbf{F}_1(\Delta)=\left.\frac{\partial}{\partial\alpha}
      e^{(\mathbf{Q}+\alpha\bm{\Gamma})\Delta}\right|_{\alpha=0},
  \qquad
  2\,\mathbf{F}_2(\Delta)=\left.\frac{\partial^{2}}{\partial\alpha^{2}}
      e^{(\mathbf{Q}+\alpha\bm{\Gamma})\Delta}\right|_{\alpha=0},
  \label{eq:tilted-derivatives}
\end{equation}
with $\bm{\Gamma}=\mathrm{diag}(\bm{\gamma})$ as above, the $(i\to j)$ entry of the first is
$\mathbb{E}[A_{0\to\Delta}\,\mathbf{1}\{X_\Delta=j\}\mid X_0=i]$ and of the second the same
expectation with $A^2_{0\to\Delta}$ in place of $A_{0\to\Delta}$; the factor of two is the two
triangular halves of the double integral the second derivative unfolds into. Dividing by the
transition probability conditions on the endpoints and dividing by $\Delta$ turns the
integral into an average, so
\begin{equation}
  \bar{\gamma}_{i\to j}
    = \frac{[\mathbf{F}_1(\Delta)]_{i\to j}}{\Delta\,P_{i\to j}(\Delta)},
  \qquad
  \bar{v}_{i\to j}
    = \frac{2[\mathbf{F}_2(\Delta)]_{i\to j}}{\Delta^2\,P_{i\to j}(\Delta)}
      - \bar{\gamma}^{\,2}_{i\to j}.
  \label{eq:gammabar-vbar}
\end{equation}

Both derivatives are evaluated from an eigendecomposition of the generator, which is the
route the implementation takes and the one the P2X2 analysis of \citet{moffatt2025bayesian}
was run on. Writing $\mathbf{Q}=\mathbf{V}\bm{\Lambda}\mathbf{W}$, with
$\bm{\Lambda}=\mathrm{diag}(\lambda_1,\dots,\lambda_K)$ the eigenvalues, $\mathbf{V}$ the
matrix of right eigenvectors and $\mathbf{W}=\mathbf{V}^{-1}$, the transition matrix is
$\mathbf{P}(\Delta)=\mathbf{V}e^{\bm{\Lambda}\Delta}\mathbf{W}$ and the two derivatives are
the same sandwich with the conductance carried in the eigenbasis. With
$\widehat{\bm{\Gamma}}=\mathbf{W}\bm{\Gamma}\mathbf{V}$ and $\circ$ the entrywise product,
\begin{align}
  \frac{\mathbf{F}_1(\Delta)}{\Delta}
    &= \mathbf{V}\bigl(\widehat{\bm{\Gamma}}\circ\mathbf{E}_2\bigr)\mathbf{W},
    &\quad (\mathbf{E}_2)_{ab} &= E_2(\lambda_a\Delta,\lambda_b\Delta),
    \label{eq:f1-spectral}\\
  \frac{2\,\mathbf{F}_2(\Delta)}{\Delta^{2}}
    &= 2\,\mathbf{V}\mathbf{M}\mathbf{W},
    &\quad M_{ab} &= \sum_{c}\widehat{\Gamma}_{ac}\,\widehat{\Gamma}_{cb}\,
                     E_3(\lambda_a\Delta,\lambda_c\Delta,\lambda_b\Delta),
    \label{eq:f2-spectral}
\end{align}
where $E_2$ and $E_3$ are the first and second divided differences of the exponential,
\begin{equation}
  E_2(x,y)=\frac{e^{x}-e^{y}}{x-y},
  \qquad
  E_3(x,y,z)=\frac{e^{x}}{(x-y)(x-z)}+\frac{e^{y}}{(y-x)(y-z)}+\frac{e^{z}}{(z-x)(z-y)}.
  \label{eq:divided-differences}
\end{equation}
Both are symmetric in their arguments and both are smooth where the denominators vanish, so
repeated and near-repeated eigenvalues are read off the coincidence limits $E_2(x,x)=e^{x}$
and $E_3(x,x,x)=e^{x}/2$, with the one-pair-coincident form between them (Supplementary
Information).
Eq.~\ref{eq:f1-spectral} is the form in which the first moment is given in the supplement of
\citet{moffatt2025bayesian}; the second moment was not stated there. The eigendecomposition
is computed once per agonist concentration and reused by every window at that concentration,
after which each window costs the $O(K^3)$ contraction of Eq.~\ref{eq:f2-spectral}.
% src: legacy/CLI_function_table.h:209-233 (Calc_eigen wrapped in Single_Thread_Memoizer, keyed on
% Agonist_concentration, so the decomposition really is shared across windows and not recomputed)
% src: legacy/qmodel.h:1681-1780 (calc_Qdt_eig, the production path; dispatch at 3147-3166 with the
% series fallback below), 814-865 (Ee and E3 with their coincidence branches), 1066-1092 (calc_eigen:
% W is tr(VL), the left eigenvectors normalized biorthogonally against V rather than V inverted, so
% W = V^-1 holds by construction; kappa_V is computed here). Read 2026-08-06.
% NAMING, and it is a real clash: what this manuscript calls E_2 is `Ee` in the code, which reserves
% `E2` for exp[0,x,y]. E_2 follows the published supplement (Moffatt & Pierdominici-Sottile 2025,
% Eq. S1), which the reader can check; E_3 is the same object under both names. Do not "align" the
% manuscript to the code here without changing the published-form claim above it.
% VERIFIED: E_2 = exp[x,y] and E_3 = exp[x,y,z] give exactly the code's kernels, and the
% normalizations match Eq. gammabar-vbar: the code's gtotal_ij is F_1/Delta and its gtotal_sqr_ij is
% 2F_2/Delta^2, because Ee/E3 are evaluated on the SCALED spectrum lambda*Delta, which carries the
% 1/Delta and 1/Delta^2. Checked at lambda->0: E_3(0,0,0)=1/2 against the direct double integral
% Delta^2/2.
% ROUTE CHANGED 2026-08-06 (Luciano). What stood here was the Van Loan block construction: the 2K
% and 3K augmented generators, F_1 and F_2 read off their upper blocks. That is correct mathematics
% and it is NOT what the code runs, so a replication following the body text would build a different
% numerical object in exactly the two places the section then discusses (repeated eigenvalues, and
% the 0/0 at unreachable transitions), and the shrinkage weight below, which is keyed to the
% conditioning of V, would have no referent at all. The block form is kept, demoted to the
% equivalent-evaluation sentence below, with its citation intact.

An equivalent evaluation avoids diagonalizing $\mathbf{Q}$ altogether. $\mathbf{F}_1$ and
$\mathbf{F}_2$ are the strictly upper blocks of the exponentials of the augmented generators
\begin{equation}
  \mathcal{Q}_2=\begin{pmatrix}\mathbf{Q} & \bm{\Gamma}\\ \bm{0} & \mathbf{Q}\end{pmatrix},
  \qquad
  \mathcal{Q}_3=\begin{pmatrix}\mathbf{Q} & \bm{\Gamma} & \bm{0}\\
                               \bm{0} & \mathbf{Q} & \bm{\Gamma}\\
                               \bm{0} & \bm{0} & \mathbf{Q}\end{pmatrix},
  \label{eq:augmented}
\end{equation}
of sizes $2K\times 2K$ and $3K\times 3K$, whose exponentials are block triangular with
$e^{\mathbf{Q}\Delta}$ on the diagonal \citep{vanloan1978computing}, so that
$e^{\mathcal{Q}_3\Delta}$ on its own returns the transition matrix and both moments. It
identifies the two objects with no spectral machinery, and where no reliable eigensolver is
at hand it is the construction to use. The route taken here is the spectral one, and the
reason is the derivatives rather than the exponential. This paper propagates the derivative
of the log-likelihood with respect to the kinetic parameters through every object the
filter builds, and the eigenvalues and eigenvectors of $\mathbf{Q}$ move with those
parameters by a first-order perturbation of the same decomposition, in closed form, while a
route built on $e^{\mathcal{Q}_3\Delta}$ carries that derivative payload through an object
of side $3K$ whose nine blocks hold three distinct matrices. Every number reported in this
paper comes from the spectral route.
% src: theory/macroir/notes/Gmean_ij_gvarij/gmean_ij_gvar_by_blocks.tex (the block derivation and
% W = 2F_2) ; legacy/parameters_derivative.h:1449-1507 (the analytic eigen derivative the last
% sentence refers to: dlambda = diag(VL^T dQ VR) and dVR = VR*Omega with Omega_jk = G_jk/(l_k - l_j),
% guarded at |l_k - l_j| > 100*sqrt(eps), plus the zero-mode gauge for the generator at 1509-1514;
% Derivative<Eigs> is memoized alongside the primitive at CLI_function_table.h:229-233).
% THIS IS SETTLED IN THE REPOSITORY AND THE PARAGRAPH FOLLOWS IT. The augmented route is not
% hypothetical here: calc_Qdt_schur (legacy/qmodel.h:2496-2542) is exactly it, ONE 3N x 3N Pade
% scaling-and-squaring on [[Q,G,0],[0,Q,G],[0,0,Q]]*dt, harvesting P, A and B from the (0,0), (0,1)
% and (0,2) blocks. It is live, and the micro path calls it (legacy/micro_full.h:1738, 1973). Its
% derivative branch is abandoned IN THE CODE and says why, at qmodel.h:2502-2510: it delegates to
% calc_Qdt_taylor because "every Derivative matmul on the 3N x 3N augmented matrix is 27x more
% expensive than the same operation on N x N (cost scales as N^3 x (1+2p))", a ~10x speedup at
% N = 51 and p = 6 with no accuracy loss. That is the sentence the body paraphrases.
% THE CAVEAT THAT KEEPS 27 OUT OF THE TEXT, and it is the author's own:
% theory/scientific-software/notes/Qdt_schur_vs_taylor_rationale.md argues the 27x is dense against
% sparse and NOT algorithmic. Only 3 of the 9 blocks are distinct; a sparse-aware version storing
% (P, F, B) costs 6 N x N matmuls per doubling, identical to calc_Qdt_taylor, which that note calls
% "the natural sparse form of the Van Loan algorithm". So the body must not claim the Van Loan
% identity is intrinsically expensive: what is expensive is carrying a derivative payload through
% the dense object, which is what the sentence says and all it says.
% NOTE for the record: the eigensolver fallback, calc_Qdt_taylor (qmodel.h:2271-2299), is itself
% that sparse Van Loan form: P(2t) = P^2, F(2t) = PF + FP, B(2t) = PB + BP + FF by doubling. No
% reported run took it; if one does, it is a Methods sentence, not this one.

Both expressions are $0/0$ where a transition is unreachable, which is not a numerical edge case
here: at zero agonist $k_{\mathrm{on}}[A]$ vanishes identically, so $P_{C\to O}(\Delta)=0$ exactly
over the whole pre-pulse segment. Conditioning on a transit that occupies a vanishing fraction of
the window leaves the chain at its two endpoint states for essentially equal times, so the limit is
the endpoint average,
\begin{equation}
  \bar{\gamma}_{i\to j}\to\tfrac12(\gamma_{i}+\gamma_{j}),
  \qquad
  \bar{v}_{i\to j}\to\Bigl(\tfrac{\gamma_{i}-\gamma_{j}}{2}\Bigr)^{2},
  \label{eq:endpoint-limit}
\end{equation}
which depends on the conductance vector alone. The implementation does not branch on the
singular case: it reads the two quotients of Eq.~\ref{eq:gammabar-vbar} as posterior means
and shrinks them towards this limit with a pseudo-count whose weight is the floating-point
noise of the spectral reconstruction, so the limit is approached smoothly and no fitted
constant enters (Methods).
% MOVED TO METHODS 2026-08-06 (Luciano). The pseudo-count itself, its two prior means, the weight
% eps_mach*kappa_F(V) and the two design choices behind them are now the first of the numerical
% safeguards at 06_methods.tex, beside the trust coefficient and the binomial floor, which is where
% a reimplementer looks for all of it at once. What stays here is the hint, because the sentence
% this replaces claimed the code TESTS for the singular case, which it does not, and leaving the
% limit unqualified in the body would put that claim back.
% src for the Methods paragraph, kept here too since this is where the limit is stated:
% legacy/qmodel.h:1691-1694 and 1743-1757 (min_P_prior = eps_mach * kappa_V), 1160-1179
% (gmean_ij_prior and gsqr_ij_prior are the endpoint averages), 1189-1202 (calc_g_ij_bayes, the
% pseudo-count quotient), 1750 (gvar_ij by subtraction, so its prior is exactly Eq. endpoint-limit).
% GAP CLOSED 2026-08-04, and it was the one that stopped a replication attempt in four of its six
% stages. The manuscript previously said only that these two objects "are computed from Q and the
% recording filter (Methods)", and neither Methods nor Appendix 1 gave a formula, a method or a
% citation, while these two symbols carry the whole av=1 against av=2 distinction, the MR/VR/IR
% algebra, the G and H matrices of Appendix 1 and the MR-IR identity.
% Sources transplanted, both verified against their own derivations:
%   theory/macroir/notes/Gmean_ij_gvarij/gmean_ij_gvar_by_blocks.tex (the two blocks, the reads,
%     W = 2F_2, the two normalisations)
%   theory/macroir/notes/Gmean_ij_gvarij/bayesian_prior_regularization_of_Qdt.md (the endpoint
%     average as the small-P limit, and why the smooth 1/sqrt(P^2+eps^2) replacement is worse:
%     it smuggles in a zero-conductance prior without flagging it)
% What stays out of the body and belongs in the Supplementary Information: the derivation of why
% F_1 and F_2 are those derivatives, the spectral evaluation of the two of them at index level with
% the coincidence branches written out, and the argument for the endpoint average. Recipe here,
% derivation there.
% AMENDED 2026-08-06 with the route change above. The pseudo-count paragraph now states the
% regularizer the code applies, which the 2026-08-04 transplant left out: it took only the
% "endpoint average as the small-P limit" half of bayesian_prior_regularization_of_Qdt.md and
% presented the limit as if the code tested for the singular case, which it does not. The rejected
% 1/sqrt(P^2+eps^2) alternative is still the note's, still out of the body.
The start-conditioned mean conductance is the $K$-vector
\begin{equation}
  \bar{\gamma}_{i\to\cdot} \;=\; \sum_{j} P_{i\to j}(\Delta)\,\bar{\gamma}_{i\to j},
  \label{eq:gammabar}
\end{equation}
which conditions on the state at the interval's start alone and is the object \texttt{INR},
\texttt{MR} and \texttt{VR} use. The boundary-conditioned mean conductance is the same
average retained on both endpoints, so it stays a $K\times K$ object rather than collapsing
to a vector, and is the object \texttt{IR} uses.
% DE-JARGONED 2026-08-06 (Luciano), here and at the two sites below. The text said "the object used
% at av=1" and "at av=2". The averaging flag is defined once in the manuscript, at
% 06_methods.tex:201, which in eLife order the reader reaches after the Results, so in the body it
% was an undefined symbol standing for a flag of the program. Naming the members instead says the
% same thing, is checkable against Table 1, and costs no notation. The flag stays in Methods, where
% it is the subject.

\paragraph{One generalization, and the whole family is inside it.}
Eq.~\ref{eq:macror} takes a channel's conductance to be a number once its state is known. That is
true of the current at an instant and of nothing else. One step of generality covers every member
below: alongside the conductance, supply a \emph{per-state variance} $\mathbf{w}$, the variance the
recorded value still has once the state is fixed. Channels being independent given their states,
those variances add over the ensemble, so
\begin{equation}
  \mathbb{E}\bigl[\bar y \bigm| \mathbf{N}\bigr]=\bm{\gamma}^\top\mathbf{N},
  \qquad
  \mathrm{Var}\bigl[\bar y \bigm| \mathbf{N}\bigr]
    =\epsilon^2_{0\to\Delta}+\textstyle\sum_s N_s\,w_s ,
  \label{eq:emission}
\end{equation}
and the extra variance is linear in the counts, exactly as the mean is. Taking the expectation over
the state distribution turns the predictive variance of Eq.~\ref{eq:macror} into
\begin{equation}
  \sigma^2 \;=\; \epsilon^2_{0\to\Delta}
    \;+\; N_{\mathrm{ch}}\,\bm{\gamma}^\top\bm{\Sigma}\bm{\gamma}
    \;+\; N_{\mathrm{ch}}\,\bm{\mu}^\top\mathbf{w},
  \label{eq:macror-w}
\end{equation}
with every other line of Eq.~\ref{eq:macror} unchanged. Setting $\mathbf{w}=\bm{0}$ returns the
published filter. This is the same device \citet{munch2022bayesian} add to the same filter to model
open-channel noise, from a different motivation; here the state-dependent term will carry the
gating variance the chosen state leaves unexplained rather than a property of the pore.

\paragraph{Four contractions, and the members differ in three of them.}
Written this way, Eq.~\ref{eq:macror} contracts a member's objects against its own state
distribution and covariance, and it does so four times:
\[
  \underbrace{\bm{\mu}^\top\bm{\gamma}}_{\text{the mean}}
  \qquad
  \underbrace{\bm{\mu}^\top\mathbf{w}}_{\text{the state-dependent noise}}
  \qquad
  \underbrace{\bm{\gamma}^\top\bm{\Sigma}\bm{\gamma}}_{\text{what the state contributes}}
  \qquad
  \underbrace{\bm{\Sigma}\bm{\gamma}}_{\text{state against current},\ \mathbf{g}}
\]
two against $\bm{\mu}$ and two against $\bm{\Sigma}$, and everything else in Eq.~\ref{eq:macror} is
a scalar. Extending the filter to a recorded value that spans the window therefore cannot require a
new algorithm. It requires choosing a state variable and supplying $(\bm{\gamma},\mathbf{w})$ on
it, and the members below are those choices. The per-state variance is not a free third choice
beside the other two: Eq.~\ref{eq:emission} defines it as whatever variance the chosen state leaves
unexplained, so naming the state names it.
% CORRECTED here relative to 02_theory.tex, which called these "the two contractions" at the end of
% the gain paragraph above. Two is the count against Sigma alone, and it undercounts even Eq. macror
% as published, where mu^T gamma is a third. With the per-state variance in the equation there are
% four. The miscount was load-bearing: it is what makes the interval variance look like an axis of
% its own, which is the reading that admits VR as a fourth member rather than as a mismatch. See the
% caption of Table 1.

\subsection*{The three steps that build the interval update}

What stops Eq.~\ref{eq:macror} from being used on the interval average is that the average is
not a function of the occupancy at any one time, so no conductance vector on the $K$ states
delivers its mean. On the $K^2$ boundary pairs one does: a channel that starts in $i$ and
ends in $j$ contributes $\bar{\gamma}_{i\to j}$ on average, so the interval average
is linear in the boundary counts again. That is the whole of the move, and three steps finish
it. Run Eq.~\ref{eq:macror} on the pairs, with $\bar{\gamma}_{i\to j}$ in place of
$\bm{\gamma}$ and the boundary covariance in place of $\bm{\Sigma}$. Charge what the
conductance still varies by once the pair is fixed to $\epsilon^2$, as though the recording
carried an extra noise of variance $\bar{v}_{i\to j}$ per channel, of a size that
depends on the pair that channel took. Marginalize over the start index, which the next window
has no use for. Every equation in this subsection is one of those three steps carried out, and
the two objects the interval likelihood is usually written with are those two contractions
read on the larger space.

The boundary state's own Gaussian follows from the prior and the propagator. Let
$\mathbf{X}_{0\to\Delta}$ be the indicator of the pair a single channel takes, one on the
$(i,j)$ it starts and ends in and zero on the other $K^2-1$, and let
$(\bm{\mu}_0,\bm{\Sigma}_0)$ be the channel's state at the window's start. It leaves $i$ and
lands in $j$ with probability $P_{i\to j}(\Delta)$, so
\begin{align}
  \mu_{i\to j} &= \mu_{0,i}\,P_{i\to j}(\Delta), \label{eq:mu-bnd}\\
  (\bm{\Sigma}_{0\to\Delta})_{(i\to j),(k\to l)} &=
    P_{i\to j}(\Delta)\,\bigl(\bm{\Sigma}_0-\mathrm{diag}(\bm{\mu}_0)\bigr)_{i k}\,
    P_{k\to l}(\Delta)
    \;+\; \bigl[\mathrm{diag}(\bm{\mu}_{0\to\Delta})\bigr]_{(i,j),(k,l)}.
  \label{eq:sig-bnd}
\end{align}
The two terms of Eq.~\ref{eq:sig-bnd} are the two sources of spread. The first carries the
uncertainty already in the prior across the window, once on each index. The second is the
uncertainty a channel known to be leaving $i$ still has about where it lands, the covariance
of a single draw from row $i$ of $\mathbf{P}(\Delta)$, which is diagonal on the pairs
because a channel takes exactly one of them. Subtracting $\mathrm{diag}(\bm{\mu}_0)$ in the
first term is what keeps the prior's own single-channel part from being counted twice, since
the second term regenerates it at the resolution of the pairs.

Summing the start index out of Eqs.~\ref{eq:mu-bnd} and~\ref{eq:sig-bnd} gives the occupancy
Gaussian at the end of the window,
\begin{align}
  \bm{\mu}_\Delta &= \mathbf{P}(\Delta)^\top \bm{\mu}_0, \label{eq:mu-prop}\\
  \bm{\Sigma}_\Delta &= \mathbf{P}(\Delta)^\top\bigl(\bm{\Sigma}_0-\mathrm{diag}(\bm{\mu}_0)\bigr)\mathbf{P}(\Delta) + \mathrm{diag}(\bm{\mu}_\Delta), \label{eq:sig-prop}
\end{align}
term by term, the diagonal of the one becoming the diagonal of the other. This is the
propagation an unconditioned member performs, and it is Eq.~\ref{eq:sig-bnd} seen from the end
of the window alone.
% MOVED UP 2026-08-06, to just after Eq. macror. The per-channel convention used to be declared
% here, after four equations that already assume it, which is where it landed when Eq. mu-prop
% opened the subsection. It now stands where the normalisation is first used.
% ADDED 2026-08-04. The convention was used consistently throughout this section and declared
% nowhere, which is worse than a contradiction because nothing in the text looks wrong. A reader
% implementing from Eq. mu-prop and Eq. sig-prop builds Sigma_0 as a total covariance and then
% every N_ch factor downstream is off by that factor. Found by a replication audit.
% SUPERSEDED IN ITS CONTENT 2026-08-06, kept for the history. The 2026-08-04 transplant was right
% that the convention had to be declared and wrong about which convention it is. It took it from
% theory/macroir/docs/Macro_IR/macroir_derivation.tex:73-95 ("all covariance matrices in the spec
% are per-channel covariances, i.e. total covariances divided by N_ch"), a derivation document, when
% the implementation states the meaning directly at legacy/qmodel.h:867-874 and states a different
% one: Sigma is the covariance of the one-hot state indicator of a SINGLE channel. Same numbers
% while the channels are independent, different object. The rule the episode illustrates is the
% repo's own: the producer code decides, a derivation note does not. Details at the correction
% comment below Eq. sig-post.
The first line of Eq.~\ref{eq:macror}, on the boundary state, is the predictive. Its mean is
\begin{equation}
  \bar{y}^{\mathrm{pred}}_{0\to\Delta} \;=\; N_{\mathrm{ch}}\,\bm{\mu}_0^\top \bar{\bm{\gamma}}_0,
  \label{eq:ypred}
\end{equation}
since $\bm{\mu}_{0\to\Delta}$ contracted against $\bar{\gamma}$ is Eq.~\ref{eq:gammabar}
summed against the prior, and its variance has three terms,
\begin{equation}
  \sigma^2_{0\to\Delta} \;=\; \epsilon^2_{0\to\Delta}
  \;+\; N_{\mathrm{ch}}\,\widetilde{\bm{\gamma}^\top\bm{\Sigma}\bm{\gamma}}
  \;+\; N_{\mathrm{ch}}\,\bm{\mu}_0^\top\mathbf{w}.
  \label{eq:pred-var}
\end{equation}
The first is the $\epsilon^2$ of Eq.~\ref{eq:macror}, the instrumental noise now taken over the
window, which carries the acquisition averaging and decreases with $\Delta$; its parametric
form is given in Methods. The third is the noise charged in step two,
$N_{\mathrm{ch}}\sum_{i,j}\mu_{i\to j}\bar{v}_{i\to j}$, which
is what Eq.~\ref{eq:pred-var} writes with $\mathbf{w}$, the
start-conditioned interval variance, once the end index is summed. The middle term is
$\bm{\gamma}^\top\bm{\Sigma}\bm{\gamma}$ of Eq.~\ref{eq:macror} evaluated on the pairs, and the
tilde is what marks that: under it the two $\bm{\gamma}$ stand for
$\bar{\gamma}_{i\to j}$ and $\bm{\Sigma}$ for the $K^2\times K^2$ boundary covariance,
not for the $K$-vector and the $K\times K$ matrix they name everywhere else. Written out with $\bm{\Sigma}_{0\to\Delta}$
of Eq.~\ref{eq:sig-bnd} and then collapsed by substituting it,
\begin{align}
  \widetilde{\bm{\gamma}^\top\bm{\Sigma}\bm{\gamma}}
  &\;=\; \sum_{i,j}\ \sum_{k,l}
    \bar{\gamma}_{i\to j}\,
    \bigl(\bm{\Sigma}_{0\to\Delta}\bigr)_{(i,j),(k,l)}\,
    \bar{\gamma}_{k\to l} \nonumber\\[2pt]
  &\;=\; \bar{\bm{\gamma}}_0^\top\bigl(\bm{\Sigma}_0-\mathrm{diag}(\bm{\mu}_0)\bigr)\bar{\bm{\gamma}}_0
    \;+\; \bm{\mu}_0^\top\bigl(\mathbf{H}\mathbf{1}\bigr),
  \label{eq:tilde}
\end{align}
where $H_{i\to j} = \bar{\gamma}^2_{i\to j}\,P_{i\to j}(\Delta)$ and $\mathbf{1}$
is the vector of ones. The two summands are the two terms of Eq.~\ref{eq:sig-bnd} in order,
and only $K$-vectors and $K\times K$ matrices survive, so the $K^2\times K^2$ covariance is
never built. Reading the tilde as
$\bar{\bm{\gamma}}_0^\top\bm{\Sigma}_0\bar{\bm{\gamma}}_0$ keeps the first summand and drops
the second, which is the single-channel spread of the interval mean across boundary pairs and
does not vanish when the occupancy covariance does.

The gain is the same contraction taken once. Built on the pairs it is the $K^2$-vector
$\sum_{k,l}(\bm{\Sigma}_{0\to\Delta})_{(i,j),(k,l)}
\bar{\gamma}_{k\to l}$, and step three sums its start index away, leaving the
$\mathbf{g}_{0\to\Delta}$ the interval update uses,
\begin{equation}
  \mathbf{g}_{0\to\Delta} \;=\; \widetilde{\bm{\gamma}^\top\bm{\Sigma}}
  \;=\; \mathbf{P}(\Delta)^\top\bigl(\bm{\Sigma}_0-\mathrm{diag}(\bm{\mu}_0)\bigr)\bar{\bm{\gamma}}_0
  \;+\; \mathbf{G}^\top\bm{\mu}_0,
  \qquad
  G_{i\to j} \;=\; \bar{\gamma}_{i\to j}\,P_{i\to j}(\Delta),
  \label{eq:vector-tilde}
\end{equation}
the covariance between the state $\mathbf{X}_\Delta$ a channel ends the window in and the
current that same channel contributes to the window average. Over the ensemble it is
$\mathrm{Cov}(\mathbf{N}_\Delta,\bar{y})=N_{\mathrm{ch}}\,\mathbf{g}_{0\to\Delta}$, with $\mathbf{N}_\Delta$
the counts by state at the end of the window, and that is the form the update uses.
Again the two summands are the two terms of Eq.~\ref{eq:sig-bnd}:
$\mathbf{G}$ carries one power of $\bar{\gamma}$ where $\mathbf{H}$ carried two, and the
end index is left open instead of being summed against $\mathbf{1}$.

The second line of Eq.~\ref{eq:macror}, marginalized the same way, is the update, with
$\delta = \bar{y}^{\mathrm{obs}}_{0\to\Delta} - \bar{y}^{\mathrm{pred}}_{0\to\Delta}$ the innovation,
\begin{align}
  \bm{\mu}^{\mathrm{post}} &= \bm{\mu}_\Delta + \frac{\mathbf{g}_{0\to\Delta}}{\sigma^2_{0\to\Delta}}\,\delta, \label{eq:mu-post}\\
  \bm{\Sigma}^{\mathrm{post}} &= \bm{\Sigma}_\Delta - \frac{N_{\mathrm{ch}}\,\mathbf{g}_{0\to\Delta}\mathbf{g}_{0\to\Delta}^\top}{\sigma^2_{0\to\Delta}}. \label{eq:sig-post}
\end{align}
Conditioning and summing the start index commute here, because the sum is linear and the
correction is rank one on the index it acts on, so the downdate crosses it unchanged, with
$\mathbf{g}_{0\to\Delta}$ standing where the boundary gain stood. What the sum discards is the posterior correlation
between the interval current and the state the window \emph{started} in, which costs nothing
inside the window and is what the recursion pays for below.

Two of the three steps are exact and one is the approximation. Eqs.~\ref{eq:mu-bnd}
and~\ref{eq:sig-bnd} are the true moments of $\mathbf{X}_{0\to\Delta}$, and so are the three
the predictive consumes. Writing $\mathbf{N}_{0\to\Delta}=\sum_c\mathbf{X}^{(c)}_{0\to\Delta}$
for the boundary counts over the ensemble, channels are independent given the pairs each of
them took, so given $\mathbf{N}_{0\to\Delta}$ the recorded average has mean
$\sum_{i,j}\bar{\gamma}_{i\to j}N_{i\to j}$ and variance
$\epsilon^2_{0\to\Delta} + \sum_{i,j}\bar{v}_{i\to j}N_{i\to j}$
exactly, and the state-dependent noise of step two, having no covariance with the counts that
carry it, enters $\sigma^2_{0\to\Delta}$ and changes neither $\bar{y}^{\mathrm{pred}}_{0\to\Delta}$ nor
$\mathbf{g}_{0\to\Delta}$. That last point is what makes charging it to the noise term legitimate rather
than convenient. What is
approximate is the Gaussian, in the two places the box above names: the counts are not
Gaussian, which is the occupancy closure, and neither is the window average of one channel
given its pair, so replacing its law by the two numbers $\bar{\gamma}_{i\to j}$ and
$\bar{v}_{i\to j}$ is the interval-signal closure. The two together still do not make
the pair $(\mathbf{N}_{0\to\Delta},\bar{y})$ jointly Gaussian, since its conditional variance
depends on the state, so the update conditions the Gaussian that carries the exact moments.
% DE-KALMANED 2026-08-06 (Luciano), here and at 08_appendix_members.tex:62. The passage read "the
% Gaussian update of the occupancy is the scalar Kalman correction" and defined g in words only,
% with the formula reachable only from Appendix 1. Two objections, one editorial and one factual.
% The formula belongs where the symbol is introduced. And the update is Gaussian conditioning of a
% joint law whose moments are computed from interval physics; naming it a Kalman correction imports
% a measurement equation y = Hx + v that this construction never writes down and that is not
% available for an interval average. The prior-art relation (integrated-measurement augmented
% Kalman, coincident to 1e-8) is conceded in the Discussion, 05_discussion.tex:127, which is the
% right place for it.
%
% REWRITTEN 2026-08-06 (Luciano), same day, on his own derivation, told directly instead of
% presented as a set of given formulas with a tilde operator arriving from nowhere: take MacroR as
% true, apply it to the BOUNDARY state, take gmean_ij as the conductance and add N*p*gvar_ij as if
% it were a state-dependent noise, then marginalise the start index. The equations fall out. What
% this buys, beyond being his: the tilde stops being a new operator and becomes MacroR's own c'Sc
% and Sc read on the K^2 pairs (so "the general operator ... is given in Appendix 1" is gone, there
% is no operator to generalise); Sigma^bnd is finally PRINTED (Eq. sig-bnd) instead of named in
% words, which was the same defect as g's; the propagation Eqs. mu-prop/sig-prop become the
% marginal of the boundary Gaussian rather than a separate starting point; the pivot of the whole
% construction gets said in one sentence (the interval average is not a function of the occupancy
% at any one time, but IS linear in the boundary counts); and the two closures of the box above
% acquire an exact address, one per approximate step, with the other two steps exact.
% VERIFIED before writing, tmp/macroir_from_macror_boundary.py: MacroR run on the K^2 pairs and
% marginalised reproduces mu-prop, sig-prop, ypred, pred-var, tilde, vector-tilde, mu-post and
% sig-post to 9.3e-16 relative, on two priors (multinomial start and a perturbed post-update one),
% and MR built the same way on the start state reproduces the mr-vr identity to 9e-16 while its
% gain genuinely differs. Eq. sig-bnd itself checked against a Monte Carlo of 4e5 replicates of 200
% channels, worst entry 2.9 standard errors over 256 entries.
% The exactness paragraph is new and is the honest statement of where the approximation is: the
% three predictive moments are exact functions of (mu_0, Sigma_0) because channels are independent
% given their boundary pairs, and the state-dependent noise contributes nothing to the
% cross-covariance, which is what licenses charging it to r^2.
% src: theory/macroir/docs/Macro_IR/macroir_derivation.tex (boundary covariance 213-363, predictive
% variance 654-1013, vector tilde 1015-1165) and elife-template.tex:360-390 (MacroR as published,
% from which Eq. macror is the state-space-agnostic transcription).
The factor $N_{\mathrm{ch}}$ in the downdate, and its absence from the mean correction, come
from the same place. All three of $\bm{\mu}$, $\bm{\Sigma}$ and $\mathbf{g}_{0\to\Delta}$ belong to one
channel, while what is observed belongs to the ensemble, so the conditioning is done on
$\mathrm{Cov}(\mathbf{N}_\Delta,\bar{y})=N_{\mathrm{ch}}\mathbf{g}_{0\to\Delta}$ against
$\mathrm{Cov}(\mathbf{N}_\Delta)=N_{\mathrm{ch}}\bm{\Sigma}_\Delta$. The rank-one
term picks up the channel count twice and the return to one channel gives one back, leaving a
single factor; the mean correction picks it up once and gives it back, leaving none.
% CONVENTION CORRECTED 2026-08-06 (Luciano). This paragraph used to say that "Sigma is carried on
% the population scale ... whereas mu and the gain g are per-channel objects", which is false on
% both halves and was the visible symptom of the convention imported at the same time as it, four
% paragraphs above: Sigma = Cov(N)/N_ch. The implementation says otherwise, and says it in the one
% place a reimplementer would look, legacy/qmodel.h:867-874: "P_Cov is stored as the *bare centered
% covariance* of the one-hot channel-state indicator: Cov(X) = diag(P_mean) - P_mean*P_mean^T.
% Sanity check: deterministic state P = e_i => Cov(X) = 0". That is a covariance of ONE channel's
% state, a primitive, not a population covariance with a factor divided out; the two coincide
% numerically while the channels are independent, which is why the equations never caught it.
% The bad convention is the spec's, theory/macroir/docs/Macro_IR/macroir_derivation.tex:73-95 ("all
% covariance matrices in the spec are per-channel covariances, i.e. total covariances divided by
% N_ch"), and it reaches this section only through the 2026-08-04 transplant recorded above. The
% notation is Milescu's and the code's, and it is the one the P2X2 manuscript uses.
% STILL OPEN, one sentence the body does not say: the two readings part company after the first
% update, because conditioning on a shared observation correlates the channels, so what the
% recursion then carries is Cov(N)/N_ch and no longer any single channel's covariance. That is the
% same non-factorization the occupancy closure is about (line 283, citing moffatt2007estimation),
% so the honest place for it is there rather than here. Luciano's call.

\subsection*{The start-conditioned members, stated forward}

The construction above reaches \texttt{IR} by choosing the boundary pair as the state. The members
that stop at the interval's start are usually presented as that construction with the collapse
brought forward, and read that way they look like abbreviations of it. They are not, and the
difference matters for what is being compared later in this section, so each is stated here from
its own assumption instead.

\texttt{MR} takes the recorded value to depend on the occupancy where the window opens. Formally
this asserts
\begin{equation}
  \bar y_{0\to\Delta} \ \perp\ X_\Delta \ \bigm|\ X_0 ,
  \label{eq:mr-assumption}
\end{equation}
that once the state a channel started in is known, the recorded average says nothing further about
where that channel ended. Everything \texttt{MR} does follows from Eq.~\ref{eq:mr-assumption}, and
the member is complete without any object larger than the $K$ states. What it must supply is a
conductance and a per-state variance on that state, and both are conditional moments of
$A_{0\to\Delta}$ given the start alone:
\begin{equation}
  (\bar\gamma_0)_{i}=\mathbb{E}\bigl[\tfrac{1}{\Delta}A_{0\to\Delta}\bigm| X_0=i\bigr],
  \qquad
  (\sigma^2_{\bar\gamma_0})_{i}
    =\mathrm{Var}\bigl[\tfrac{1}{\Delta}A_{0\to\Delta}\bigm| X_0=i\bigr],
  \label{eq:mr-objects}
\end{equation}
the first being Eq.~\ref{eq:gammabar} and the second the per-state variance
Eq.~\ref{eq:macror-w} then asks for. Neither needs a state larger than $K$. Both come from the same
tilted semigroup of Eq.~\ref{eq:tilted-derivatives} that any interval member needs, but from its
derivatives contracted against the vector of ones rather than read entry by entry: what
\texttt{MR} requires are the row sums, and the row sums are integrals of $K$-vectors. This is worth
stating plainly because it is what makes \texttt{MR} a construction in its own right, reached
without ever forming a pair.

Eq.~\ref{eq:mr-assumption} is an assumption about the process and can be examined on its own terms.
It is false for any window of nonzero length: a channel that starts shut and delivers current
during the window is more likely to have ended open than one that delivers none, so the recorded
average does carry information about the end state beyond what the start state gives.
\texttt{IR} is the member that declines to assume it, and the boundary pair is the state on which
the recorded average is linear with no conditional expectation taken at all. \texttt{MR} is
therefore a correctly derived filter for a measurement model that is false, rather than an
approximate version of \texttt{IR}; how much the falsehood costs is Eq.~\ref{eq:cdiff} below.

\subsection*{What separates the start-conditioned members from the boundary one}

Run as the same three steps on the $K$ states instead of the pairs, with
$\bar{\bm{\gamma}}_0$ in place of $\bm{\gamma}$, two
algebraic differences separate them from the one that conditions on both ends. The first is
the form of the third term of Eq.~\ref{eq:pred-var}.
\texttt{MR} carries the total per-start-state interval variance, the spread of the
interval-averaged conductance of a channel known only to have started in $i$, while
\texttt{VR} and \texttt{IR} carry the residual variance left once the end state is also
conditioned on, $\sum_{j} P_{i\to j}(\Delta)\,\bar{v}_{i\to j}$. By the law
of total variance the two forms differ by the spread of the interval mean across end
states,
\begin{equation}
  \underbrace{\sum_{j} P_{i\to j}(\Delta)\bigl[\bar{v}_{i\to j}
    + \bigl(\bar{\gamma}_{i\to j}-\bar{\gamma}_{i\to\cdot}\bigr)^2\bigr]}_{\text{total, }\texttt{MR}}
  \;=\; \underbrace{\sum_{j} P_{i\to j}(\Delta)\,\bar{v}_{i\to j}}_{\text{residual, }\texttt{VR},\,\texttt{IR}}
  \;+\; \mathrm{Var}_{j}\bigl[\bar{\gamma}_{i\to j}\mid i\bigr],
  \label{eq:mr-vr}
\end{equation}
so \texttt{MR} charges to the observation variance a spread that the end state, once
conditioned on, would resolve. The second difference is the boundary cross-covariance term
\begin{equation}
  N_{\mathrm{ch}}\,\widetilde{\bm{\gamma}^\top\bm{\Sigma}\bm{\gamma}}
  \label{eq:mr-ir}
\end{equation}
which \texttt{IR} retains and both start-conditioned members drop. It enters twice, in the
second term of Eq.~\ref{eq:pred-var} and in the gain $\mathbf{g}_{0\to\Delta}$ of
Eqs.~\ref{eq:mu-post} and~\ref{eq:sig-post}.

The two differences are not independent, and the tempting reading of them, one step that
moves the variance followed by one that moves the gain, is wrong in both directions. Taken
from a common prior they cancel in the predicted variance: what the boundary term adds to
the second term of Eq.~\ref{eq:pred-var} is
$\bm{\mu}_0^\top\mathrm{Var}_{j}[\bar{\gamma}_{i\to j}\mid i]$, which is what
the residual form removes from the third, so from the same state \texttt{MR} and \texttt{IR}
assign the interval average the same variance by two different decompositions and the whole
of what separates them is the gain (Appendix~\ref{app:members}). In the other direction, the predicted
variance divides the gain, so changing the variance form changes the update as well:
\texttt{VR} removes that spread from the third term without restoring it in the second,
predicts less variance than either from the same state, and its posterior departs from
\texttt{MR}'s at once. The equality is therefore between two maps from a prior to a
predicted variance, and not between two recordings. Once the gains differ the priors differ,
and the variances the two members report along one recording differ with them, which is what
the Results measure across the regime grid.
% src: legacy/qmodel.h:4568-4619, read 2026-08-06, and verified three ways rather than asserted:
% (i) term by term, gSg at av=2 is gmean_i^T·SmD·gmean_i + p·(gtotal_ij o gmean_ij)·1 against
% gmean_i^T·P_Cov·gmean_i at av<=1, while ms switches total (gsqr_i - gmean_i^2) against residual
% (gsqr_i - (gtotal_ij o gmean_ij)·1); the cross term enters gSg with + and ms with - and cancels,
% and MR's contraction against the full P_Cov contributes exactly the gamma^T diag(mu) gamma the
% total form subtracts. (ii) numerically on random fields, difference 0.0, including deliberately
% inconsistent Qdt fields: the cancellation is structural. (iii) on the figure-1 dumps, where MR and
% IR agree to the last printed digit at every interval in which they still share a prior and part
% company at the first update after that. THE QUALIFICATION IS LOAD-BEARING AND WAS MISSING HERE:
% an earlier draft of this paragraph said "MR and IR predict the same variance and differ only in
% the gain" with no "from the same state", which reads as a statement about a recording, where it is
% false. papers/1_method/figures_build_plan.md:199-220 already had the identity, the numbers and the
% retirement of the "changes nothing else" slogan; program.md:78, decisions.md:161 and
% CONTINUE_HERE.md:88 forbid that slogan in four documents, and 05_discussion.tex:101 states it
% correctly ("from the same state"). This paragraph carried the retired form until now.

\subsection*{Both members against the exact law}

Neither member is measured against the other. Both are approximations to one joint law, that of
$(\mathbf{N}_\Delta,\bar y)$ under the true process, and given
$(\bm{\mu}_0,\bm{\Sigma}_0)$ that law has exact first two moments which either member can be held
against. Two elementary decompositions supply them, and they say different things about the two
blocks.

\paragraph{The variance: both members are exact, and their slots differ.}
By the law of total variance, conditioning on the counts where the window opens,
\begin{equation}
  \mathrm{Var}\bigl[\bar y\bigr]
  = \epsilon^2_{0\to\Delta}
  + \underbrace{\mathbb{E}\bigl[\mathrm{Var}(\bar y \mid \mathbf{N}_0)\bigr]}_{N_{\mathrm{ch}}\,\bm{\mu}_0^\top\mathbf{w}}
  + \underbrace{\mathrm{Var}\bigl[\mathbb{E}(\bar y \mid \mathbf{N}_0)\bigr]}_{N_{\mathrm{ch}}\,\bar{\bm{\gamma}}_{0}^\top\bm{\Sigma}_0\bar{\bm{\gamma}}_{0}} ,
  \label{eq:ltv}
\end{equation}
which is Eq.~\ref{eq:pred-var} read with the start-conditioned objects: the middle term is
$\bar{\bm{\gamma}}_0^\top\bm{\Sigma}_0\bar{\bm{\gamma}}_0$, the tilde dropped, and $\mathbf{w}$ is
the total slot of Eq.~\ref{eq:mr-vr}. Read with the boundary objects instead, the tilde of
Eq.~\ref{eq:tilde} in the middle and the residual slot in $\mathbf{w}$, the same quantity is
Eq.~\ref{eq:pred-var} again. \textbf{Both members report the exact
variance of the recorded value} given $(\bm{\mu}_0,\bm{\Sigma}_0)$, so they agree with each other
for the reason that each agrees with the truth, and nothing is measured by that agreement. This is
Eq.~\ref{eq:mr-vr} read from the outside: the same second moment sits in the state term of one
member and in the noise term of the other, which is why the totals coincide.

The variable conditioned on in Eq.~\ref{eq:ltv} is the count vector and not one channel's state.
Eq.~\ref{eq:emission} gives
$\mathrm{Var}[\bar y\mid\mathbf{N}_0]=\epsilon^2_{0\to\Delta}+\sum_s N_{0,s}w_s$ and
$\mathbb{E}[\bar y\mid\mathbf{N}_0]=\bar{\bm{\gamma}}_0^\top\mathbf{N}_0$, so the second underbrace
is $\bar{\bm{\gamma}}_0^\top\mathrm{Cov}(\mathbf{N}_0)\bar{\bm{\gamma}}_0$ and the per-channel
convention supplies the factor of $N_{\mathrm{ch}}$ in both. Conditioning on a single channel's
$X_0$ returns the same expressions only while $\bm{\Sigma}_0$ is still that channel's own
covariance, which it stops being at the first update.

\paragraph{The covariance: one member is exact and the other is short by a named amount.}
By the law of total covariance, conditioning on the same counts and carried on the per-channel
scale,
\begin{equation}
  \frac{1}{N_{\mathrm{ch}}}\mathrm{Cov}\bigl[\mathbf{N}_\Delta,\ \bar y\bigr]
  = \underbrace{\frac{1}{N_{\mathrm{ch}}}\mathrm{Cov}\bigl[\mathbb{E}(\mathbf{N}_\Delta\mid \mathbf{N}_0),\ \bar y\bigr]}_{\textstyle \mathbf{g}^{\mathtt{MR}}_{0\to\Delta}}
  + \ \frac{1}{N_{\mathrm{ch}}}\mathbb{E}\bigl[\mathrm{Cov}(\mathbf{N}_\Delta,\ \bar y \mid \mathbf{N}_0)\bigr] .
  \label{eq:ltc}
\end{equation}
The first term is the start-conditioned gain
$\mathbf{g}^{\mathtt{MR}}_{0\to\Delta}=\mathbf{P}(\Delta)^\top\bm{\Sigma}_0\bar{\bm{\gamma}}_0$ exactly. The
second is what Eq.~\ref{eq:mr-assumption} sets to zero: given the counts the channels are
conditionally independent, so it is the per-channel conditional covariance summed over the
population, and a conditional independence is precisely the statement that it vanishes. It does not
vanish. \texttt{IR} carries the whole of Eq.~\ref{eq:ltc}, so the difference between the two
members in this block is that second summand and nothing else,
\begin{equation}
  \mathbf{g}^{\mathtt{IR}}_{0\to\Delta}-\mathbf{g}^{\mathtt{MR}}_{0\to\Delta}
  \;=\;\Bigl(\textstyle\sum_{i}\mu_{i\to j}\,D_{i\to j}\Bigr)_{j},
  \qquad
  D_{i\to j}=\bar{\gamma}_{i\to j}-(\bar\gamma_0)_{i} ,
  \label{eq:cdiff}
\end{equation}
which is Eq.~\ref{eq:vector-tilde} minus $\mathbf{P}(\Delta)^\top\bm{\Sigma}_0\bar{\bm{\gamma}}_0$
written out. The deviation $D_{i\to j}$ is a property of the process and not of either
algorithm: it is how far the interval average of a channel that ended in $j$ sits from the
average over all the places it might have ended. It has zero row sums,
$\sum_{j}P_{i\to j}(\Delta)D_{i\to j}=0$, since $(\bar\gamma_0)_{i}$ is that very
average, and it therefore has zero mean under $\mu_{i\to j}$ as well, so the gap between the
two forms of Eq.~\ref{eq:mr-vr} is exactly its variance under the law of the pairs.
Eq.~\ref{eq:cdiff} does not depend on $\bm{\Sigma}_0$ at all.

\paragraph{Why the same quantity is invisible in one block and decisive in the other.}
$D$ enters the variance squared and the covariance once, and the two blocks answer to its zero row
sums differently. A second moment can sit in either of the two variance terms without the total
noticing, which is Eq.~\ref{eq:mr-vr}. A \emph{first} moment cannot: a noise term has no first
moment, so when $D$ is collapsed against the end index there is nowhere for it to go, and its
column profile, which is what carries the information about where the channel ended, is lost. That
asymmetry between $D^2$ and $D$ is the whole of the difference between the two members, and it is a
consequence of Eq.~\ref{eq:mr-assumption} rather than a defect of the arithmetic that implements
it.
% VERIFIED 2026-08-09, and these are the checks the tmp/MR_IR_derivation family carried in
% comments where no reader could see them; the visible statement is in "What was verified" below.
% Eq. cdiff: identity to 1e-17 in 36 cases, independent of Sigma_0, and against a Monte Carlo of the
% CTMC path. Three equivalent writings of the same term, all verified: weighted by mu_0 with the
% per-start variance; weighted by mu^bnd with D^2 per pair; and Var(D) under mu^bnd. They agree
% because E[D] = 0 under that law.
% src of the derivation: papers/1_method/decisions/recompute/mr_vs_ir_boxes.py

\subsection*{The same construction in the linear-filtering frame}

Everything above was Gaussian conditioning: a joint law over (state, recorded value), its three
blocks, and Bayes. Nothing was borrowed from linear filtering, and the derivation does not need it.
The correspondence with that literature is nevertheless exact, it is what places this construction
against forty years of work on the same measurement problem, and a reader arriving from there
should be able to check the identification rather than take it on faith. The relation is conceded
in the Discussion; the equations that make it checkable are here.

For a state $\mathbf{x}$ and a scalar observation $y$ jointly Gaussian, the filter's update is
\begin{equation}
  \mathbf{x}^{\mathrm{post}}=\mathbf{x}^{\mathrm{prior}}
    +\frac{\mathrm{Cov}(\mathbf{x},y)}{\mathrm{Var}(y)}\bigl(y-\mathbb{E}[y]\bigr),
  \qquad
  \mathbf{P}^{\mathrm{post}}=\mathbf{P}^{\mathrm{prior}}
    -\frac{\mathrm{Cov}(\mathbf{x},y)\,\mathrm{Cov}(\mathbf{x},y)^\top}{\mathrm{Var}(y)} ,
  \label{eq:kf}
\end{equation}
which is Eq.~\ref{eq:macror} up to the per-channel convention, the source of the factors of
$N_{\mathrm{ch}}$ there and not here. The two frames do not differ in the update at all: both are
the conditioning of a joint Gaussian, and what is usually called the gain is the ratio of the cross
block to the observation variance. The entire content of either derivation is the computation of
the three blocks $\mathbb{E}[y]$, $\mathrm{Var}(y)$ and $\mathrm{Cov}(\mathbf{x},y)$ for a recorded
value that is an average over the window.

The standard difficulty is that $\bar y$ is not a function of the state at any single time, so no
observation matrix $\mathbf{H}$ satisfies $\bar y=\mathbf{H}\mathbf{x}_\Delta+v$. The device is
state augmentation: carry the running integral
$z_t=\int_0^{t}\bm{\gamma}^\top\mathbf{N}(s)\,\mathrm{d}s$ alongside the counts, so that
$(\mathbf{N},z)$ propagates jointly and the recorded value is $z_\Delta/\Delta$ plus instrumental
noise, which \emph{is} linear in the augmented state. This is the treatment of an observation that
is a flow rather than a stock, and it is a known construction rather than a new one
\citep{zadrozny1988gaussian}. Augmentation moves the problem rather than solving it: the filter now
needs the variance of the accumulator and its covariance with the occupancy at the end of the
window, and in general those are not free. For a continuous-time Markov chain they are, and by the
two derivatives already written as Eq.~\ref{eq:tilted-derivatives}: the $(i\to j)$ entry of
$\mathbf{F}_1$ is $\mathbb{E}[A_{0\to\Delta}\mathbf{1}\{X_\Delta=j\}\mid X_0=i]$ and of
$2\mathbf{F}_2$ the same with $A^2_{0\to\Delta}$, which is exactly what the augmented filter asks for.

Condition on the counts where the window opens; given those the channels are independent, and the
three blocks are
\begin{align}
  \mathbb{E}\bigl[\bar y\bigr] &= N_{\mathrm{ch}}\,\bm{\mu}_0^\top\bar{\bm{\gamma}}_0,
  \label{eq:kf-mean}\\
  \mathrm{Var}\bigl[\bar y\bigr] &= \epsilon^2_{0\to\Delta} + \frac{\mathrm{Var}(z_\Delta)}{\Delta^2},
  \label{eq:kf-var}\\
  \frac{\mathrm{Cov}\bigl[\mathbf{N}_\Delta,\bar y\bigr]}{N_{\mathrm{ch}}}
    &= \frac{\mathrm{Cov}\bigl[\mathbf{N}_\Delta,z_\Delta\bigr]}{N_{\mathrm{ch}}\Delta}
    = \mathbf{P}(\Delta)^\top\bm{\Sigma}_0\bar{\bm{\gamma}}_0
      + \Bigl(\textstyle\sum_{i}\mu_{i\to j}D_{i\to j}\Bigr)_{j} .
  \label{eq:kf-cross}
\end{align}
Eq.~\ref{eq:kf-mean} is Eq.~\ref{eq:ypred}, Eq.~\ref{eq:kf-var} is Eq.~\ref{eq:pred-var}, and
Eq.~\ref{eq:kf-cross} is $\mathbf{g}^{\mathtt{MR}}_{0\to\Delta}$ plus Eq.~\ref{eq:cdiff}, which is to say the
$\mathbf{g}_{0\to\Delta}$ of Eq.~\ref{eq:vector-tilde}. The second term of Eq.~\ref{eq:kf-cross} needs nothing
from the boundary construction: for one channel, $\mathbf{F}_1$ divided by $\Delta$ and by the
transition probability gives
$\mathbb{E}[\mathbf{1}\{X_\Delta=j\}\,\tfrac{1}{\Delta}A_{0\to\Delta}\mid X_0=i]
= P_{i\to j}(\Delta)\bar{\gamma}_{i\to j}$, the product of the two conditional
expectations is $P_{i\to j}(\Delta)(\bar\gamma_0)_{i}$, and their difference is
$P_{i\to j}(\Delta)D_{i\to j}$. The deviation of Eq.~\ref{eq:cdiff} is therefore what the
first derivative of the tilted semigroup leaves once the product of the two first moments is taken
out, and the agreement between the two routes is a result rather than an input.

\textbf{The augmented-state filter and the boundary-state filter compute the same three blocks},
so they perform the same update and return the same posterior over the occupancy. Three things
about that follow, and they are the reason the derivation above is given in the other order.

The augmented state has $K+1$ entries and the boundary state has $K^2$, and they deliver the same
three blocks. \emph{The pair is a device of the derivation and not information the filter must
carry}, which is the same conclusion Eqs.~\ref{eq:mu-prop} and~\ref{eq:sig-prop} reached from the
other side. The per-state variance disappears in the augmented frame: given $(\mathbf{N}_\Delta,
z_\Delta)$ the recorded value is deterministic, so what Eq.~\ref{eq:pred-var} splits between its
second and third terms sits entirely inside $\mathrm{Var}(z_\Delta)$, and that is why the cross
block needs only $\mathbf{F}_1$ while the variance block needs $\mathbf{F}_2$ as well. And the
equivalence holds per window, with a reset: the augmented filter must discard the posterior on $z$
and open the next window with $z=0$, because the next window's integral is a fresh one and is
independent of the old given the occupancy.

\texttt{MR} is not the augmented filter under any choice of blocks. It is the filter for a
misspecified measurement model,
\begin{equation}
  \bar y = \mathbf{H}\mathbf{N}_0 + v(\mathbf{N}_0),
  \qquad
  \mathbf{H}=\bar{\bm{\gamma}}_0^\top,
  \qquad
  \mathrm{Var}\bigl[v\mid \mathbf{N}_0\bigr]=\epsilon^2_{0\to\Delta}+\mathbf{N}_0^\top\mathbf{w},
  \label{eq:mr-kf}
\end{equation}
whose expectation over $\mathbf{N}_0$ returns Eq.~\ref{eq:pred-var}. The measurement noise depends
on the state, which is outside the textbook filter: with that extension the two frames give the
same equations for \texttt{MR} as well, and without it \texttt{MR} has no linear-filtering form at
all.

Where the frames genuinely differ is the model rather than the arithmetic. A Kalman derivation
presumes a linear-Gaussian state process, and the occupancy counts of a channel population are not
one: they are multinomial, their transition noise depends on the state they leave, and no
linear-Gaussian system generates them. The derivation above presumes none of that. It computes the
exact first two moments of the true process, over the pairs where the recorded average is genuinely
linear, and replaces the law by a Gaussian at the two places the box above names. The two routes
arrive at the same equations because the conditioning of a joint Gaussian consumes nothing but the
first two moments, and both supply the same ones. Two consequences follow. Naming the update a
Kalman correction imports a measurement equation $\bar y=\mathbf{H}\mathbf{x}+v$ that the
construction never writes down and that is not available for an interval average without the
augmentation; and $\bm{\Sigma}$ has a meaning in one frame and a convention in the other, since
what this section carries is the covariance of one channel's state indicator and not the covariance
matrix of a linear system. The pair survives as the presentation even though the augmented state
reaches the same blocks with fewer entries, because on the pair the recorded average is linear in
the state with no conditional expectation taken, so exactness is structural rather than a separate
calculation; the collapse hands over the next window's prior at no cost; and the per-state variance
slot is what makes \texttt{R}, \texttt{MR} and \texttt{IR} one construction with three readings.
The augmentation gives no comparably natural form to the other two members.

\subsection*{What was verified, and to what tolerance}

 The identities of this section were checked
against a direct contraction over the $K^2$ pairs rather than asserted. Eqs.~\ref{eq:tilde}
and~\ref{eq:vector-tilde} reproduce the explicit contraction of the boundary covariance in $108$
random cases to $7\times10^{-18}$ and $3\times10^{-16}$; the commutation of the start-index
marginalization with the update, which is what licenses Eqs.~\ref{eq:mu-post}
and~\ref{eq:sig-post} being written on the $K$ states, holds in $36$ cases to $2\times10^{-16}$;
Eq.~\ref{eq:cdiff} holds to $1\times10^{-17}$ in $36$ cases, independently of $\bm{\Sigma}_0$, and
against a Monte Carlo of the continuous-time path; Eqs.~\ref{eq:mu-bnd} and~\ref{eq:sig-bnd} were
checked against a Monte Carlo of $4\times10^{5}$ replicates of $200$ channels, worst entry $2.9$
standard errors over $256$ entries; and the two forms of Eq.~\ref{eq:mr-vr} agree with the forms
that go through the pair-resolved variance to $1\times10^{-15}$ in $16$ cases. The whole
construction, run on the pairs and marginalized, reproduces
Eqs.~\ref{eq:mu-prop},~\ref{eq:sig-prop},~\ref{eq:ypred},~\ref{eq:pred-var},~\ref{eq:tilde},~\ref{eq:vector-tilde},~\ref{eq:mu-post}
and~\ref{eq:sig-post} to $9.3\times10^{-16}$ relative on two priors, a multinomial start and a
perturbed post-update one.
% src: papers/1_method/decisions/recompute/mr_vs_ir_boxes.py and tmp/macroir_from_macror_boundary.py.
% These numbers were carried in LaTeX comments in every previous version of this material, where no
% reader could see them. For a section whose purpose is that someone can reproduce the construction,
% the checks belong in the text.

\end{appendixbox}
